\subsection{\texorpdfstring{Preesquemas regulares y propiedad $(R_k)$. Criterio de normalidad de Serre.}{Preesquemas regulares y propiedad (R\string_k). Criterio de normalidad de Serre}}

\begin{env}[5.8.1]
\label{IV.2.s5.8.1-fr}
Recordemos $(\mathbf{0}_{\mathrm I},4.1.4)$ que un espacio anillado
$(X,\mathcal O_X)$ se dice \emph{regular} en un punto $x\in X$ si
$\mathcal O_x$ es un anillo \emph{regular}. Cuando se trate de preesquemas, en este
capítulo solo utilizaremos esta terminología cuando $X$ sea \emph{localmente
noetheriano}.
\end{env}

\begin{definition}[5.8.2]
\label{IV.2.s5.8.2-fr}
Se dice que un preesquema localmente noetheriano $X$ es \textup{regular en codimensión
$\leq k$}, o que \textup{posee la propiedad $(R_k)$}, si el conjunto de los puntos en
los que $X$ no es regular es de codimensión $>k$ \textup{(dicho de otro modo (5.1.3),}
si, para todo $x\in X$, la relación $\dim(\mathcal O_x)\leq k$ implica que
$\mathcal O_x$ es regular).
\end{definition}

Decir que $X$ es regular significa que $X$ posee la propiedad $(R_k)$ para
todo $\mathrm{k}$.

Si $X=\operatorname{Spec}(A)$, donde $A$ es un anillo \emph{noetheriano}, se dirá
que $A$ posee la propiedad $(R_k)$ si $X$ posee esta propiedad; decir que $X$
es regular significa que el anillo $A$ es \emph{regular} (\textbf 0, 17.3.6). Para un
preesquema localmente noetheriano cualquiera $X$, se dirá que $X$ \emph{posee la
propiedad $(R_k)$ en un punto} $x\in X$ si el anillo local $\mathcal O_x$
posee la propiedad $(R_k)$; esto quiere decir, pues, que para toda
\emph{generización} $x'$ de $x$ en $X$, la relación
$\dim(\mathcal O_{x'})\leq k$ implica que $\mathcal O_{x'}$ es un anillo local
regular. Decir que $X$ es regular en un punto $x$ equivale a decir que $X$
verifica la propiedad $(R_n)$ para \emph{todo} $n\geq0$ en el punto $x$, en virtud
de (\textbf 0, 17.3.6).

\begin{proposition}[5.8.3]
\label{IV.2.s5.8.3-fr}
Si $k$ es un cuerpo y $X$ un preesquema localmente de tipo finito sobre $k$,
entonces, para todo $x\in X$, existe una vecindad abierta de $x$ en $X$
isomorfa a un subesquema de un $k$-esquema regular.
\end{proposition}

En efecto, existe una vecindad abierta afín $U$ de $x$ isomorfa a un
$k$-esquema de la forma $\operatorname{Spec}(A)$, donde $A$ es una
$k$-álgebra de tipo finito; por consiguiente, $A$ es isomorfa a un cociente
de un álgebra de polinomios $k[T_1,\ldots,T_n]$, y se sabe que esta última
es un anillo \emph{regular} (\textbf 0, 17.3.7).

\begin{env}[5.8.4]
\label{IV.2.s5.8.4-fr}
En virtud de (5.8.2), decir que $X$ posee la propiedad $(R_0)$ significa que,
para todo punto maximal $x$ de $X$, el anillo $\mathcal O_x$ es un \emph{cuerpo}
(\textbf 0, 17.1.4), dicho de otro modo, que $X$ es \emph{reducido} en ese punto. Como el
conjunto $U$ de los $x\in X$ en los que $X$ es reducido es \emph{abierto} (pues el

\oldpage[IV-2]{108}

nilradical de $X$ es un $\mathcal O_X$-módulo coherente (\textbf I, 6.1.1 y
$\mathbf{0}_{\mathrm I}$, 5.2.2)), es equivalente decir que $X$ posee la propiedad
$(R_0)$ o que el conjunto $U$ es \emph{denso en todas partes}. Por consiguiente:
\end{env}

